\documentclass[11pt,a4paper]{article}
\usepackage[italian]{babel}
\usepackage{fontspec}
\defaultfontfeatures{Ligatures=TeX}
\IfFontExistsTF{TeX Gyre Pagella}{\setmainfont{TeX Gyre Pagella}}{\setmainfont{DejaVu Serif}}
\IfFontExistsTF{TeX Gyre Heros}{\setsansfont{TeX Gyre Heros}}{\setsansfont{DejaVu Sans}}
\IfFontExistsTF{Latin Modern Mono}{\setmonofont{Latin Modern Mono}}{\setmonofont{DejaVu Sans Mono}}
\IfFontExistsTF{TeX Gyre Schola}{\newfontfamily\handfont{TeX Gyre Schola}}{\newfontfamily\handfont{DejaVu Serif}}
\usepackage[a4paper,margin=18mm]{geometry}
\usepackage{microtype}
\usepackage{xcolor}
\usepackage{tikz}
\usetikzlibrary{calc,positioning,decorations.pathmorphing,arrows.meta,patterns,shapes.geometric}
\usepackage[most]{tcolorbox}
\usepackage{tabularx,array,booktabs,longtable}
\usepackage{enumitem}
\usepackage{fancyhdr}
\usepackage{graphicx}
\usepackage{ragged2e}
\usepackage{setspace}
\usepackage{hyperref}
\hypersetup{hidelinks,pdftitle={Notti di Yharnam - Dossier degli Handout},pdfauthor={Materiale amatoriale per uso privato}}

\definecolor{ink}{HTML}{191716}
\definecolor{blood}{HTML}{5A171B}
\definecolor{ash}{HTML}{69635E}
\definecolor{paper}{HTML}{F3EFE5}
\definecolor{oldpaper}{HTML}{E8DECA}
\definecolor{church}{HTML}{E7E4DC}
\definecolor{blueprint}{HTML}{E8ECEB}
\definecolor{rulegray}{HTML}{B8B0A4}

\pagestyle{fancy}
\fancyhf{}
\renewcommand{\headrulewidth}{0pt}
\renewcommand{\footrulewidth}{0pt}
\fancyfoot[C]{\scriptsize\color{ash}Materiale di gioco - non distribuire l'indice ai giocatori}

\newcommand{\gmcode}[1]{\begin{tikzpicture}[remember picture,overlay]\node[anchor=south east,font=\tiny\sffamily,color=ash] at ($(current page.south east)+(-6mm,5mm)$) {#1};\end{tikzpicture}}
\newcommand{\ornament}{\par\smallskip\begin{center}\begin{tikzpicture}\draw[blood,line width=.7pt] (-2.7,0)--(-.35,0);\draw[blood,line width=.7pt] (.35,0)--(2.7,0);\node[text=blood] at (0,0) {$\diamond$};\end{tikzpicture}\end{center}\smallskip}
\newcommand{\handtitle}[2]{\begin{center}{\Large\bfseries\textsc{#1}}\\[2mm]{\small\itshape #2}\end{center}\ornament}
\newcommand{\churchseal}{\begin{tikzpicture}[baseline=-.5ex,scale=.32]\draw[blood,line width=1pt] (0,0) circle (1);\draw[blood,line width=.8pt] (0,-.72)--(0,.75);\draw[blood,line width=.8pt] (-.55,.2)--(.55,.2);\draw[blood,line width=.6pt] (-.38,-.3)--(0,-.68)--(.38,-.3);\end{tikzpicture}}
\newcommand{\huntermark}{\begin{tikzpicture}[baseline=-.5ex,scale=.32]\draw[ink,line width=1pt] (-.8,.8)--(0,-.8)--(.8,.8);\draw[ink,line width=1pt] (-.55,.25)--(.55,.25);\draw[ink,line width=1pt] (-.35,-.15)--(.35,-.15);\end{tikzpicture}}

\newtcolorbox{gmnote}[1][]{colback=paper,colframe=blood,boxrule=.6pt,arc=0mm,title=#1,fonttitle=\bfseries\sffamily,left=3mm,right=3mm,top=2mm,bottom=2mm}
\newtcolorbox{documentbox}[1][]{colback=oldpaper,colframe=ink,boxrule=.6pt,arc=0mm,left=7mm,right=7mm,top=7mm,bottom=7mm,#1}
\newtcolorbox{churchbox}[1][]{colback=church,colframe=blood,boxrule=1pt,arc=0mm,left=8mm,right=8mm,top=7mm,bottom=7mm,borderline={.4pt}{2pt}{blood},#1}
\newtcolorbox{bluebox}[1][]{colback=blueprint,colframe=ash,boxrule=.7pt,arc=0mm,left=6mm,right=6mm,top=5mm,bottom=5mm,#1}

\setlength{\parindent}{0pt}
\setlength{\parskip}{5pt}
\newcommand{\newhandout}[3]{\clearpage\thispagestyle{empty}\gmcode{#1}\handtitle{#2}{#3}}

\begin{document}

% ----------------------------------------------------------------
% INDICE GM
% ----------------------------------------------------------------
\thispagestyle{fancy}
\begin{center}
{\Huge\bfseries\textsc{Dossier degli Handout}}\\[2mm]
{\large Fate: Notti di Yharnam - Il Progetto}
\end{center}
\ornament

\begin{gmnote}[Uso del fascicolo]
Questa prima pagina è riservata al GM. Ogni documento successivo occupa una pagina autonoma e può essere stampato, separato e consegnato senza spiegazioni. Il codice Hxx compare soltanto in piccolo nell'angolo inferiore destro. Gli handout fondamentali non sostituiscono le informazioni automatiche previste dalla Guida del GM: le rendono tangibili e confrontabili.
\end{gmnote}

\scriptsize
\renewcommand{\arraystretch}{1.05}
\begin{longtable}{>{\bfseries}p{8mm}p{39mm}p{31mm}p{72mm}}
\toprule
ID & Documento & Nodo / consegna & Funzione investigativa \\
\midrule
H01 & Taccuino di Elias I & Casa Ward & Stabilisce che Elias indagava su pesi, turni e registri. \\
H02 & Cartina annotata & Casa Ward & Offre piste simultanee: Fonderia, Archivio, Ospedale. \\
H03 & Frammento di progetto & Studio / Officina & Collega i componenti moderni alle vecchie invenzioni di Elias. \\
H04 & Registro della Fonderia & Fonderia & Mostra incongruenze tra peso dichiarato e verificato. \\
H05 & Ricevuta non registrata & Fonderia / magazzino & Prova l'esistenza di una cassa fuori registro diretta al Laboratorio Alchemico. \\
H06 & Pagina dell'Archivio & Archivio Ecclesiastico & Collega ordini frammentati, Ospedale e Laboratorio; introduce il sigillo interno. \\
H07 & Ordine interno della Chiesa & Archivio / Padre Aldren & Rivela il trasferimento riservato senza spiegare il vero Progetto. \\
H08 & Referto medico di Elias & Ospedale Vecchio & Conferma sintomi iniziali ma lucidità conservata. \\
H09 & Nota privata del dottor Sile & Ospedale & Mostra paura, denuncia e organismo incompleto. \\
H10 & Avviso di caccia & Rifugio dei Cacciatori & Spiega perché Elias è considerato preda intelligente e pericolosa. \\
H11 & Annotazioni di Miriam Voss & Biblioteca Proibita & Fornisce il vocabolario dell'imitazione del cordone e dell'“ascolto”. \\
H12 & Taccuino di Elias II & Biblioteca / Laboratorio & Dimostra che il Laboratorio Alchemico è un transito e che Elias è sceso sotto terra. \\
H13 & Inventario alchemico & Laboratorio Alchemico & Fa emergere materiali assenti, condotte e trasferimenti sotterranei. \\
H14 & Schema delle condotte & Laboratorio / fogne & Permette ai giocatori di pianificare l'accesso al nodo finale. \\
H15 & Tavola dell'apparato & Il Progetto & Mostra l'uso distorto dei progetti di Elias nel supporto biologico. \\
H16 & Protocollo sperimentale & Il Progetto & Rivela finalità, tessuti non consensuali e rischio di attivazione. \\
H17 & Ultimi appunti di Elias & Finale & Chiarisce il sabotaggio e lascia aperta la decisione morale. \\
\bottomrule
\end{longtable}
\normalsize

% ----------------------------------------------------------------
% H01
% ----------------------------------------------------------------
\newhandout{H01}{Appunti sparsi}{Foglio strappato da un taccuino da lavoro}
\begin{documentbox}
\handfont
\large
\textbf{Turno della notte - terza campana}

Cassa 17: dichiarati \textbf{412 libbre}.\\
Bilancia orientale: \textbf{391}.\\
Bilancia piccola: \textbf{391 e un quarto}.

Non è un errore. Due strumenti non sbagliano nello stesso modo.

\medskip
Valvola di compensazione, seconda serie. Mia.\\
O una copia abbastanza buona da esserlo.

\medskip
Perché proprio il turno di notte?\\
Perché il registro viene ricopiato dopo la chiusura?\\
Chi decide quali casse non devono arrivare?

\vfill
\hfill\textit{Controllare l'Archivio. Non parlarne con Marta finché non so.}
\end{documentbox}

% ----------------------------------------------------------------
% H02
% ----------------------------------------------------------------
\newhandout{H02}{Pianta piegata della città}{Annotazioni a grafite attribuite a Elias Ward}
\begin{bluebox}
\centering
\begin{tikzpicture}[x=1cm,y=1cm,font=\small]
\draw[ash,line width=.4pt] (-7,-5) rectangle (7,5);
\draw[ash!60,line width=.4pt] plot[smooth cycle,tension=.6] coordinates {(-6,1)(-5,4)(-1,4.5)(2.5,4)(5.8,2)(6,-2)(3,-4.3)(-2,-4.5)(-5.5,-2.5)};
\draw[ash!55,line width=.5pt] (-5.5,-1)--(-3,1)--(-1,2.5)--(2,2)--(5,3);
\draw[ash!55,line width=.5pt] (-4,3)--(-2,1)--(0,-1)--(3,-2)--(5,-1);
\draw[ash!55,line width=.5pt] (-5,-3)--(-2,-2)--(0,0)--(2,2.5)--(4,4);
\draw[ash!55,line width=.5pt] (-1,4)--(-.5,1)--(-1,-2)--(0,-4);

\node[draw=ink,fill=paper,inner sep=2pt] (f) at (-4.6,-1.8) {Fonderia};
\node[draw=ink,fill=paper,inner sep=2pt] (a) at (1.2,2.9) {Archivio};
\node[draw=ink,fill=paper,inner sep=2pt] (o) at (3.9,.1) {Ospedale};
\node[draw=ink,fill=paper,inner sep=2pt] (l) at (1.8,-2.8) {Lab. Alchemico};
\node[draw=ash,fill=paper,inner sep=2pt] at (-1.5,1.0) {Mercato};
\node[draw=ash,fill=paper,inner sep=2pt] at (-.7,-3.5) {Fogne};
\node[draw=ash,fill=paper,inner sep=2pt] at (4.2,3.7) {Cattedrale};

\draw[blood,line width=1.1pt] (f) circle (.95);
\draw[blood,line width=1.1pt] (a) circle (.85);
\draw[blood,line width=1.1pt] (o) circle (.85);
\draw[blood,dashed,line width=.8pt,->] (f)--(l);
\draw[blood,dashed,line width=.8pt,->] (l)--(o);
\draw[blood,dotted,line width=.8pt] (o)--(a);
\node[rotate=-8,text=blood,font=\itshape] at (-2.8,-3.3) {non segue la strada};
\node[rotate=6,text=blood,font=\itshape] at (3.5,-1.5) {sotto?};
\node[rotate=-5,text=ink,font=\itshape] at (-2.2,3.5) {registro ricopiato qui};
\end{tikzpicture}
\end{bluebox}

% ----------------------------------------------------------------
% H03
% ----------------------------------------------------------------
\newhandout{H03}{Regolatore di pressione - Serie II}{Frammento recuperato dal vecchio studio di Elias}
\begin{bluebox}
\begin{tikzpicture}[x=1cm,y=1cm,font=\footnotesize]
\draw[ink,line width=.7pt] (-5,0) rectangle (-2,2.2);
\draw[ink,line width=.7pt] (-2,1.1)--(0,1.1);
\draw[ink,line width=.7pt] (0,.3) rectangle (2.3,1.9);
\draw[ink,line width=.7pt] (2.3,1.1)--(4.8,1.1);
\draw[ink,line width=.7pt] (-3.5,0)--(-3.5,-2.1);
\draw[ink,line width=.7pt] (-4.4,-1.3) rectangle (-2.6,-2.1);
\draw[ink,line width=.7pt] (.4,.3)--(.4,-1.7);
\draw[ink,line width=.7pt] (.4,-1.7)--(3.5,-1.7);
\draw[ink,line width=.7pt] (3.5,-2.2) rectangle (4.8,-1.2);
\draw[ink,->] (-5.8,1.1)--(-5,1.1);
\draw[ink,->] (4.8,1.1)--(5.7,1.1);
\draw[ink,decorate,decoration={zigzag,segment length=4pt,amplitude=2pt}] (-1.5,.5)--(-.5,1.7);
\node at (-3.5,1.1) {camera primaria};
\node at (1.15,1.1) {compensatore};
\node at (-3.5,-1.65) {scarico};
\node at (4.15,-1.7) {ritorno};
\draw[ash,<->] (-5,2.7)--(-2,2.7) node[midway,fill=blueprint] {18 pollici};
\draw[ash,<->] (0,2.4)--(2.3,2.4) node[midway,fill=blueprint] {11};
\node[align=left,anchor=west] at (-5,-3) {\textit{Nota originaria:} progettato per mantenere stabile\\la pressione durante trasfusioni prolungate.};
\node[align=left,anchor=west,text=blood] at (-5,-4) {\textit{Annotazione successiva:} se collegato in serie,\\può sostenere un carico molto più grande di un corpo umano.};
\end{tikzpicture}
\vfill
\small\itshape Firma parziale: E. W---. Il margine destro è annerito dal fuoco.
\end{bluebox}

% ----------------------------------------------------------------
% H04
% ----------------------------------------------------------------
\newhandout{H04}{Fonderia Municipale}{Registro di pesatura - turno notturno}
\begin{documentbox}
\sffamily
\begin{center}\Large\bfseries REGISTRO DI PESATURA E USCITA\end{center}
\vspace{2mm}
\renewcommand{\arraystretch}{1.45}
\begin{tabularx}{\linewidth}{c c c c X}
\toprule
Data & Cassa & Dichiarato & Verificato & Destinazione \\
\midrule
12 & 14 & 286 lb & 286 lb & Ospedale Vecchio \\
12 & 15 & 174 lb & 175 lb & Laboratorio Alchemico \\
13 & 16 & 322 lb & 322 lb & Ospedale Vecchio \\
13 & 17 & 412 lb & \textbf{391 lb} & \textit{non indicata} \\
14 & 18 & 208 lb & 208 lb & Laboratorio Alchemico \\
14 & 19 & 146 lb & 146 lb & Ospedale Vecchio \\
15 & 20 & 297 lb & 296 lb & Laboratorio Alchemico \\
\bottomrule
\end{tabularx}

\vspace{7mm}
\textbf{Osservazioni del turno}

Cassa 17 sigillata prima della verifica. Ordine di non aprire. La differenza è stata confermata da due bilance. Il sigillo non appartiene all'amministrazione municipale.

\vfill
\begin{tabularx}{\linewidth}{X X}
Caporeparto: \textit{Marta Ghal} & Verificatore: \textit{E. Ward}\\
\end{tabularx}
\end{documentbox}

% ----------------------------------------------------------------
% H05
% ----------------------------------------------------------------
\newhandout{H05}{Ricevuta di trasferimento}{Copia carbone priva di registrazione ufficiale}
\begin{churchbox}
\begin{center}
\churchseal\quad {\Large\bfseries SERVIZIO RISERVATO DI TRASPORTO}\quad\churchseal
\end{center}
\ornament
\begin{tabularx}{\linewidth}{>{\bfseries}p{42mm}X}
Numero collo & 17-B \\
Origine & Fonderia Municipale, cancello orientale \\
Destinazione dichiarata & Laboratorio Alchemico \\
Destinazione effettiva & \rule{70mm}{.4pt} \\
Peso accettato & 391 libbre e un quarto \\
Contenuto & Camera di compensazione, vetriatura, minuteria d'argento \\
Consegna & Dopo la terza campana; ingresso inferiore \\
Autorizzazione & Sigillo del \textit{Servizio dell'Ascolto} \\
\end{tabularx}

\vspace{8mm}
\textit{Nota a margine, altra mano:} «Non portare nell'ala pubblica. La pressione deve restare costante durante la discesa.»

\vfill
\hfill Firma del vetturale: \textit{T. Harrow}
\end{churchbox}

% ----------------------------------------------------------------
% H06
% ----------------------------------------------------------------
\newhandout{H06}{Archivio Ecclesiastico}{Pagina 44 - commesse straordinarie}
\begin{churchbox}
\begin{center}{\large\bfseries CHIESA DELLA CURA - REGISTRO DELLE COMMESSE}\end{center}
\small
\renewcommand{\arraystretch}{1.35}
\begin{tabularx}{\linewidth}{c X X c}
\toprule
Rif. & Materiale & Destinatario ufficiale & Sigillo \\
\midrule
44-11 & Tubazioni di rame, giunti sterili & Ospedale Vecchio & Cappellania \\
44-12 & Vetriatura spessa, supporti molleggiati & Laboratorio Alchemico & \churchseal \\
44-13 & Tessuti di sutura, sali conservanti & Ospedale Vecchio & Infermeria \\
44-14 & Regolatori di pressione, serie II & Laboratorio Alchemico & \textbf{\churchseal\,\churchseal} \\
44-15 & Lastre d'argento e fili sottili & \textit{voce raschiata} & \textbf{\churchseal\,\churchseal} \\
44-16 & Carbone purificato, campane minute & Torre Astronomica & Accademia \\
\bottomrule
\end{tabularx}

\vspace{5mm}
Le voci con doppio sigillo sono sottratte alla consultazione ordinaria. I documenti accessori devono essere trasferiti all'ala inferiore dell'Ospedale e distrutti dopo la verifica.

\vfill
\textit{In matita, nel margine:} «Il sigillo non compare nei registri degli ordini. Chi lo riconosce?»
\end{churchbox}

% ----------------------------------------------------------------
% H07
% ----------------------------------------------------------------
\newhandout{H07}{Ordine di servizio interno}{Copia incompleta, contrassegnata “non archiviare”}
\begin{churchbox}
\begin{center}
\churchseal\quad {\Large\bfseries AD USO DEL CLERO AUTORIZZATO}\quad\churchseal
\end{center}
\vspace{5mm}
Ai responsabili dell'Ospedale Vecchio, del Laboratorio Alchemico e dell'Archivio:

Le apparecchiature comprese nelle commesse 44-12, 44-14 e 44-15 dovranno essere trasferite attraverso i passaggi di servizio senza contatto con il personale ordinario. Ogni richiesta di chiarimento sarà ricondotta al programma di miglioramento emoterapico.

Il soggetto maschile identificato come \textbf{Elias Ward}, già ingegnere e ora operaio municipale, ha posto domande relative alle consegne. Non deve essere arrestato pubblicamente. Qualora si presenti in una struttura della Chiesa, trattenerlo con cortesia e informare Padre Aldren Vale.

Le sue carte devono essere recuperate integre.

\vspace{8mm}
\textit{La presente disposizione non costituisce ammissione dell'esistenza di alcun programma ulteriore.}

\vfill
\hfill Per ordine del Servizio dell'Ascolto
\end{churchbox}

% ----------------------------------------------------------------
% H08
% ----------------------------------------------------------------
\newhandout{H08}{Ospedale Vecchio}{Referto di visita non protocollato}
\begin{documentbox}
\sffamily
\begin{center}{\Large\bfseries SCHEDA DI OSSERVAZIONE CLINICA}\end{center}
\renewcommand{\arraystretch}{1.25}
\begin{tabularx}{\linewidth}{>{\bfseries}p{38mm}X}
Paziente & Elias Ward, anni 42 circa \\
Professione & Operaio di fonderia; precedenti mansioni tecniche \\
Motivo della visita & Insonnia, irritabilità, ipersensibilità olfattiva, cefalea \\
Temperatura & Elevata ma stabile \\
Pupille & Reattive; lieve dilatazione \\
Linguaggio & Coerente, rapido, ossessivamente preciso \\
Memoria & Integra \\
Segni bestiali & Nessuna alterazione ossea; unghie ispessite; risposta aggressiva a odore di sangue \\
\end{tabularx}

\vspace{5mm}
\textbf{Osservazioni:} il paziente comprende pienamente il proprio stato e teme una progressione. Ha rifiutato emoterapia aggiuntiva. Ha chiesto se «una macchina possa trasmettere una febbre senza toccare il malato».

\textbf{Valutazione:} non classificabile come bestia al momento della visita. Sorveglianza consigliata.

\vfill
Medico: \textit{Dott. Oren Sile}\hfill Data: nove giorni fa
\end{documentbox}

% ----------------------------------------------------------------
% H09
% ----------------------------------------------------------------
\newhandout{H09}{Nota mai spedita}{Carta personale del dottor Oren Sile}
\begin{documentbox}
\handfont
Caro collega,

non so più a chi affidare queste righe. Ward non delirava. Parlava troppo in fretta, ma ogni numero che mi mostrò coincideva con ciò che ho visto nell'ala inferiore.

Tre notti fa mi ordinarono di preparare sei pazienti per un «trasferimento terapeutico». Nessuno di loro aveva dato consenso. Nel vano di servizio vidi una vasca coperta, collegata a pompe industriali. Qualcosa al suo interno si mosse quando le campane della cappella tacquero.

Ho informato Padre Aldren della visita di Ward. Credevo di proteggerlo. Da allora l'ala è chiusa e due pazienti sono scomparsi dai registri.

Se Ward torna, non consegnarlo senza garanzie. Se non torna, temo di sapere dove sia andato.

\vfill
\hfill Oren
\end{documentbox}

% ----------------------------------------------------------------
% H10
% ----------------------------------------------------------------
\newhandout{H10}{Avviso di caccia provvisorio}{Affisso all'interno del Rifugio dei Cacciatori}
\begin{documentbox}
\begin{center}
\huntermark\quad {\Large\bfseries SOGGETTO DA INTERCETTARE}\quad\huntermark
\end{center}
\vspace{4mm}
\begin{tabularx}{\linewidth}{>{\bfseries}p{38mm}X}
Nome & Elias Ward \\
Età & circa quarantadue anni \\
Corporatura & media, spalle incurvate, mano destra ustionata \\
Ultimo abbigliamento & cappotto da operaio, borsa di attrezzi, pistola corta \\
Sintomi riferiti & insonnia, aggressività, olfatto alterato, forza intermittente \\
Ultima traccia & area compresa tra Ospedale Vecchio, Laboratorio Alchemico e condotti fognari \\
\end{tabularx}

\vspace{5mm}
Il soggetto conserva intelligenza eccezionale e conoscenza di macchinari, serrature, pressione e dispositivi di sicurezza. Una trasformazione completa potrebbe produrre una preda capace di preparare trappole e sabotare infrastrutture cittadine.

\textbf{Ordine:} contenere se possibile. Abbattere al primo segno di perdita irreversibile del controllo. Recuperare taccuini e strumenti.

\vfill
\hfill Capitana Roland Kreiss
\end{documentbox}

% ----------------------------------------------------------------
% H11
% ----------------------------------------------------------------
\newhandout{H11}{Marginalia sull'Ascolto}{Annotazioni attribuite alla dottoressa Miriam Voss}
\begin{documentbox}
\handfont
\textbf{Estratto copiato da un testo pthumeriano:}

\begin{quote}\itshape
Il figlio non chiama con la bocca. Il legame ascolta per lui. Dove il legame è reciso, il sangue ricorda la strada ma non il nome.
\end{quote}

\textbf{Note:}

- Il «legame» non sembra un organo nel senso comune. È funzione, passaggio, accordatura.

- Un cordone autentico non può essere fabbricato. Si potrebbe però imitare temporaneamente la sua funzione con tre elementi: \textit{materia vivente}, \textit{sangue affine}, \textit{ritmo stabile}.

- Le pompe della Chiesa possono fornire il ritmo. Le campane minute potrebbero sostituire ciò che i testi chiamano «voce senza suono».

- L'apparato non creerebbe un essere superiore. Creerebbe qualcosa che tenta di essere udito.

\vspace{6mm}
\textit{Annotazione verticale nel margine:} «Se Ward ha riconosciuto il ritmo, allora il suo progetto è già parte dell'esperimento.»
\end{documentbox}

% ----------------------------------------------------------------
% H12
% ----------------------------------------------------------------
\newhandout{H12}{Appunti sparsi - seconda pagina}{Carta macchiata d'olio e piegata più volte}
\begin{documentbox}
\handfont
\large
Il Laboratorio Alchemico non consuma abbastanza carbone per tutte le macchine che riceve.

Non produce abbastanza scarico.\\
Non conserva abbastanza personale.\\
Non è il laboratorio.

\medskip
Le pompe lavorano a valle. Si sente il ritorno nei tubi quando l'edificio è fermo.

Tre accessi possibili:

1. condotto dell'Ospedale;\\
2. ascensore merci sotto il deposito;\\
3. scarico verso le fogne, ma il livello sale alla quarta campana.

\medskip
Ho lasciato la valvola orientale instabile. Tre colpi, pausa, due colpi: seguire la vibrazione.

\vfill
\hfill\textit{Non portare Eleanor. Non ancora.}
\end{documentbox}

% ----------------------------------------------------------------
% H13
% ----------------------------------------------------------------
\newhandout{H13}{Laboratorio Alchemico}{Inventario di reparto - beni in transito}
\begin{documentbox}
\sffamily
\begin{center}{\Large\bfseries INVENTARIO DELLA SETTIMANA}\end{center}
\renewcommand{\arraystretch}{1.35}
\begin{tabularx}{\linewidth}{X c c X}
\toprule
Voce & Entrata & Presente & Note \\
\midrule
Regolatori Serie II & 4 & 0 & trasferiti inferiormente \\
Vetriature grandi & 6 & 1 & cinque trasferite \\
Fili d'argento & 18 bobine & 2 & sedici trasferite \\
Sali conservanti & 12 casse & 3 & nove trasferite \\
Carbone purificato & 20 sacchi & 18 & consumo ordinario \\
Campane minute & 24 & 0 & non registrare in uscita \\
Tessuto biologico sigillato & 7 contenitori & 0 & consegna diretta \\
\bottomrule
\end{tabularx}

\vspace{7mm}
\textbf{Trasferimento inferiore:} usare l'ascensore soltanto dopo la terza campana. In caso di guasto, aprire lo scarico orientale e attendere il segnale: tre colpi, pausa, due colpi.

\vfill
\textit{Correzione a matita:} «Valvola orientale sabotata. Chiamare un tecnico che non faccia domande.»
\end{documentbox}

% ----------------------------------------------------------------
% H14
% ----------------------------------------------------------------
\newhandout{H14}{Schema delle condotte inferiori}{Copia tecnica non aggiornata}
\begin{bluebox}
\centering
\begin{tikzpicture}[x=1cm,y=1cm,font=\footnotesize,>=Latex]
\node[draw,align=center,minimum width=3cm,minimum height=1cm] (h) at (0,3.5) {Ospedale Vecchio\\ala inferiore};
\node[draw,align=center,minimum width=3cm,minimum height=1cm] (a) at (-4,0) {Laboratorio\\Alchemico};
\node[draw,align=center,minimum width=2.7cm,minimum height=1cm] (s) at (4,0) {Scarico\\fognario};
\node[draw=blood,align=center,line width=1pt,minimum width=3.2cm,minimum height=1.2cm] (p) at (0,-3.3) {CAMERA INFERIORE\\accesso riservato};
\node[draw,diamond,aspect=2] (v) at (0,1.3) {V-Est};
\node[draw,circle] (lift) at (-2,-1.7) {A};
\node[draw,circle] (gate) at (2,-1.7) {G};

\draw[->,line width=1pt] (h)--(v);
\draw[->,line width=1pt] (a)--(v);
\draw[->,line width=1pt] (s)--(gate);
\draw[->,line width=1pt] (v)--(p);
\draw[->,line width=1pt] (a)--(lift)--(p);
\draw[->,line width=1pt] (gate)--(p);
\draw[blood,dashed,line width=1pt] (-.8,1.3) arc (180:350:.8);
\node[text=blood,align=left] at (3.7,2.7) {V-Est instabile\\ritmo 3--2};
\node[align=left] at (-4,-3.3) {A: ascensore merci\\chiave doppia};
\node[align=left] at (4,-3.3) {G: grata di scarico\\alta marea alla IV campana};
\end{tikzpicture}

\vfill
\small\textit{Sul retro, una nota di Elias:} «La camera è sotto l'Ospedale, ma il controllo arriva dal Laboratorio. Tagliare una sola linea non basta.»
\end{bluebox}

% ----------------------------------------------------------------
% H15
% ----------------------------------------------------------------
\newhandout{H15}{Apparato di risonanza ematica}{Tavola tecnica affissa nella camera inferiore}
\begin{bluebox}
\begin{tikzpicture}[x=1cm,y=1cm,font=\footnotesize]
\draw[ink,line width=.8pt] (-1.5,-1.4) ellipse (2.1 and 2.7);
\draw[ink,line width=.7pt] (-1.5,1.3) ellipse (1.2 and .55);
\draw[ink,line width=.7pt] (-1.5,-3.5) ellipse (1.5 and .4);
\draw[ink,line width=1pt] (-4.8,-.2) rectangle (-3,1.7);
\draw[ink,line width=1pt] (2.2,-.2) rectangle (4,1.7);
\draw[ink,line width=.7pt] (-3,1)--(-2.7,1)--(-2.7,.3)--(-2.1,.3);
\draw[ink,line width=.7pt] (2.1,.3)--(2.7,.3)--(2.7,1)--(2.2,1);
\draw[ink,line width=.7pt] (-3.9,-.2)--(-3.9,-2.5)--(-2.3,-2.5);
\draw[ink,line width=.7pt] (3.1,-.2)--(3.1,-2.5)--(-.7,-2.5);
\draw[ink,decorate,decoration={coil,segment length=4pt,amplitude=3pt}] (-.6,.8)--(-.6,-1.7);
\draw[ink,decorate,decoration={coil,segment length=4pt,amplitude=3pt}] (-2.4,.8)--(-2.4,-1.7);
\node at (-3.9,.75) {pompa A};
\node at (3.1,.75) {pompa B};
\node[align=center] at (-1.5,-.3) {supporto\\biologico};
\node[align=left,anchor=west] at (4.4,2.1) {1. stabilizzare pressione};
\node[align=left,anchor=west] at (4.4,1.5) {2. introdurre sangue affine};
\node[align=left,anchor=west] at (4.4,.9) {3. avviare campane};
\node[align=left,anchor=west,text=blood] at (4.4,.1) {4. attendere risposta};
\draw[ash,<->] (-1.5,2.8)--(-1.5,1.85) node[midway,right] {ritmo};
\draw[ash,<->] (-5.4,-.2)--(-5.4,1.7) node[midway,left] {Serie II};
\node[align=left,anchor=west] at (-5.2,-4.2) {\textit{Correzione a mano di Elias:} Le pompe non tengono in vita ciò che è nella vasca.\\Tengono aperto qualcosa attraverso di essa.};
\end{tikzpicture}
\end{bluebox}

% ----------------------------------------------------------------
% H16
% ----------------------------------------------------------------
\newhandout{H16}{Protocollo VII - Prima risonanza}{Documento del Servizio dell'Ascolto}
\begin{churchbox}
\begin{center}
\churchseal\quad {\Large\bfseries PROTOCOLLO VII}\quad\churchseal\\[2mm]
{\small Prima risonanza dell'imitazione ombelicale}
\end{center}
\vspace{4mm}
\begin{enumerate}[leftmargin=7mm]
\item Stabilizzare il supporto biologico con due pompe Serie II in opposizione.
\item Introdurre sangue affine in tre incrementi. Il terzo deve coincidere con la settima sequenza di campane.
\item Collegare i fili d'argento ai tessuti auricolari e spinali.
\item Se il supporto reagisce alla presenza di persone non sottoposte a emoterapia, annotare nomi e sintomi.
\item Non interrompere il ritmo dopo l'inizio della risposta. Una cessazione improvvisa può trasferire lo stress biologico al soggetto affine più vicino.
\end{enumerate}

\textbf{Materiale biologico autorizzato:} campioni ospedalieri 3, 4 e 6; residui pthumeriani; tessuti neonatali non reclamati.

\textbf{Esito atteso:} non nascita, ma apertura temporanea di un canale percettivo.

\textbf{Esito avverso:} agitazione delle bestie, perdita di coscienza, crescita incontrollata, risposta esterna non identificata.

\vfill
Autorizzazione: \textit{Madre Celestine Vey}\hfill Non duplicare
\end{churchbox}

% ----------------------------------------------------------------
% H17
% ----------------------------------------------------------------
\newhandout{H17}{Ultimi appunti di Elias}{Foglio trovato nei vani tecnici del Progetto}
\begin{documentbox}
\handfont
\large
Le valvole orientali sono aperte al contrario. L'apparato può ancora funzionare, ma non a piena pressione.

Ho separato il ritorno dell'Ospedale. Se attivano adesso, perderanno il supporto oppure trasferiranno il carico su chiunque abbia sangue compatibile nelle vicinanze.

Non so quale delle due cose sia peggiore.

\medskip
La vasca non contiene un figlio. Contiene un tentativo di chiamarne uno.

Le mie macchine danno il ritmo. Il sangue dà la strada. Qualcosa dall'altra parte deve soltanto decidere se ascoltare.

\medskip
La febbre peggiora. Ho ferito un uomo e ricordo di aver scelto di non ucciderlo. Questo significa che sono ancora io, oppure significa soltanto che la bestia sa aspettare?

\medskip
Se Eleanor è qui, ditele che non sono rimasto per la conoscenza.

\textit{Ho continuato a studiare perché non sapevo come distruggerlo senza uccidere tutti quelli collegati.}

\vfill
\hfill E.
\end{documentbox}

\end{document}
